\section{Du terrain au système d’information : Domaine d'étude et acteur}
Avant la présentation des modèles conceptuels, il est nécessaire de définir le domaine d'étude et d'identifier les acteurs. En Merise, un acteur est une personne, entité, organisation, dispositif ou système externe échangeant des informations avec le domaine, il peut être interne ou externe. Dans notre cas, ces acteurs ont été identifiés à partir des rôles réellement observés et exercés lors de l'expérimentation, permettant ainsi d'établir leur correspondance avec le système d’information.

\begin{longtable}{|p{7cm}|p{3.8cm}|p{2.2cm}|}
    \caption{Acteurs du système d'information, avec leur type et leur rôle sur le terrain}
    \label{tab:acteurs_si} \\
    \hline
    \textbf{Rôle observé sur le terrain (chapitre 3 et protocole)} & \textbf{Acteur du système d'information} & \textbf{Type} \\
    \hline
    \endfirsthead
    \hline
    \textbf{Rôle observé sur le terrain (chapitre 3 et protocole)} & \textbf{Acteur du système d'information} & \textbf{Type} \\
    \hline
    \endhead
    Équipe expérimentale, relevant les données du sol à l'aide du capteur 8-en-1 deux fois par jour (matin et soir), calculant $ET_0$/$ET_c$, décidant d'irriguer ou non, appliquant les dosages d'amendement (protocole) & Technicien & Interne \\
    \hline
    Équipe expérimentale, finançant l'achat du capteur 8-en-1 et des composants de l'amendement (biochar, coque d'arachide, compost, biofertilisant) & Bailleur & Interne \\
    \hline
    Responsable scientifique & Utilisateur & Externe \\
    \hline
    Le capteur 8-en-1, mesurant humidité, température, pH, conductivité électrique, N, P, K & Capteur 8-en-1 & Interne \\
    \hline
    L'application météo du smartphone (AccuWeather) utilisée pendant l'expérimentation, remplacée dans l'application par l'API OpenWeather & AccuWeather & Externe \\
    \hline
    Le fournisseur du capteur 8-en-1 & Fournisseur capteur & Externe \\
    \hline
    Le fournisseur des composants de l'amendement (charbon de bois, coque d'arachide, compost, engrais bio) & Fournisseur intrant & Externe \\
    \hline
\end{longtable}

Cette correspondance entre l'expérience vécue et les acteurs modélisés assure que le système conçu automatise une pratique réelle documentée et testée, basée sur les conditions agronomiques et climatiques observées sur le champ expérimental Agro TIC de l'USSEIN.
